\documentclass[../main.tex]{subfiles}
\begin{document}
\section*{\centering 概要}
\begin{tcolorbox}[
  title=執筆メモ,
  colback=yellow!8,
  colframe=black!30,
  breakable
]
\begin{itemize}
  \item 物理目標
  \begin{itemize}
    \item ダイクォーク相関
    \item ハドロン内部の構造
  \end{itemize}
  \item 計画されている実験と本研究のコントリビューション
  \begin{itemize}
    \item チャームバリオンの分光実験(J-PARC E50)
    \item 連続読み出し式DAQを使うが、非常に高いデータレートが想定されており、オンラインフィルタを用いてデータ量を削減する必要がある
    \item 本番で使う想定のオンラインフィルターを実装し、テスト実験データを用いて性能評価を行った
    \item 性能評価の結果を用いて本番で必要とされる計算資源を見積もった
  \end{itemize}
  \item 本研究の内容
  \begin{itemize}
    \item 3種のハイレベルフィルターを実装
    \item テスト実験T103のデータを用いて性能評価
    \item 本番で必要とされる計算資源を見積もった
  \end{itemize}
  \item 本研究の結果
\end{itemize}
\end{tcolorbox}

\begin{tcolorbox}[
  title=過去の修論の切り口,
  colback=yellow!8,
  colframe=black!30,
  breakable
]
\begin{itemize}
  \item 林
  \begin{itemize}
    \item (アブストがしっかりしている)
    \item ハドロンの基底状態や低励起状態がCQMで記述できるが、それを超えると説明がつかないことがいくつもある
    \item それに対してダイクォーク相関を有効自由度とした記述が試みられているが、それは実験的に確認されていない
    \item チャームなら見れる(と考えられている)
    \item 
  \end{itemize}
  \item 山本
  \begin{itemize}
    \item (Missing Math法というありえない誤字をしてる)
    \item 計画している実験
    \item 荷電K,piを取りたいからドリフトチェンバーが必要で、その読み出し基板を開発
    \item $\upbeta$線とSP8のビーム試験で性能評価
  \end{itemize}
  \item 
\end{itemize}
\end{tcolorbox}


\end{document}
